\documentclass[11pt,a4paper]{article}
\usepackage[T1]{fontenc}
\usepackage[utf8]{inputenc}
\usepackage[main=italian,provide=*]{babel}
\usepackage{amsmath,amssymb}
\usepackage{geometry}
\geometry{margin=2.5cm}
\usepackage{booktabs}
\usepackage{longtable}
\usepackage{seqsplit}
\usepackage{array}
\usepackage{hyperref}
\hypersetup{colorlinks=true,linkcolor=blue,citecolor=blue,urlcolor=blue}

\title{Convenzioni sito$\leftrightarrow$qubit nel progetto — documento di riferimento}
\author{Note di lavoro — Tesi triennale, Samuele Schiavi\\ Relatore: Prof.\ Alessandro Chiesa}
\date{}

\begin{document}
\maketitle

\section{Perché questo documento}

Il progetto usa \textbf{due convenzioni diverse e incompatibili} per
mappare i siti fisici (spin $1,2,3,\dots$) sugli indici dei qubit Qiskit,
a seconda del modulo. Nessuna delle due è "sbagliata" — ogni modulo è
internamente autoconsistente, con i propri self-test a precisione
macchina — ma \textbf{mescolare stati o operatori provenienti da moduli
con convenzioni diverse senza tradurli produce errori di segno
silenziosi}, difficili da diagnosticare perché non danno crash: danno
numeri plausibili ma sbagliati (tipicamente invisibili su osservabili
reali come $\langle\sigma^z\rangle$, ma che flippano il segno su
$\langle\sigma^y\rangle$ o su correlatori misti).

Questo documento è organizzato \textbf{nello stesso ordine in cui il
lavoro di tesi è stato svolto}: dimero (N=2) prima, poi trimero ad anello
(N=3, catena chiusa), poi trimero a catena aperta (N=3). Per ciascun
blocco tematico: quali file lo compongono, a cosa serve ciascuno, quale
famiglia di convenzione usa e \emph{perché} è stata scelta proprio quella
in quel punto del lavoro — non per riportare una regola statica, ma per
ricostruire la logica (spesso storica, non sempre "ottimale a
tavolino") con cui le scelte si sono succedute.

\section{La regola generale di Qiskit (little-endian)}

Qiskit indicizza i qubit in modo che, in una stringa di Pauli o in un ket
scritto come stringa di bit a $n$ caratteri, \textbf{il carattere più a
sinistra corrisponde al qubit con indice più alto} ($n-1$), quello più a
destra al qubit $0$. In simboli: la posizione $k$ della stringa (contando
da sinistra, $k=0,\dots,n-1$) agisce sul qubit $n-1-k$. Questa regola è
fissa (built-in in Qiskit); quello che \emph{è} una scelta del progetto è
\emph{quale sito fisico si mette a quale posizione della stringa}.

Due famiglie di scelta ricorrono nel progetto:

\begin{itemize}
\item \textbf{Famiglia 1} — "sito 1 a sinistra": sito$_i\to$qubit$(N{-}i)$.
Per $N=2$: sito1$\to$qubit1, sito2$\to$qubit0. Per $N=3$: sito1$\to$qubit2,
sito2$\to$qubit1, sito3$\to$qubit0. Vantaggio: la label Pauli letta da
sinistra a destra coincide con l'ordine fisico $|q_1q_2\dots q_N\rangle$
usato in tutta la derivazione teorica — nessuna reindicizzazione fra
teoria e codice.
\item \textbf{Famiglia 2} — "sito 1 a destra": sito$_i\to$qubit$(i{-}1)$.
Per $N=2$: sito1$\to$qubit0, sito2$\to$qubit1.
\end{itemize}

\section{Parte 1 — Dimero (N=2)}

\subsection{1.1 Struttura statica: benchmark esatto, VQE}

\textbf{Perché Famiglia 1.} È il punto di partenza dell'intero progetto:
\texttt{dimer\_exact.py} fissa la convenzione apposta per far leggere la
label Pauli nello stesso ordine della notazione fisica
$|q_1q_2\rangle$ — una scelta di leggibilità, non imposta da Qiskit, fatta
una volta all'inizio e poi ereditata da tutto ciò che dipende
direttamente da questo file.

\begin{longtable}{@{}p{4.7cm}p{6.3cm}p{2.5cm}@{}}
\toprule
File & A cosa serve & Mappa\\
\midrule
\endhead
\texttt{\seqsplit{dimer\_exact.py}} & Benchmark esatto (diagonalizzazione) del dimero, autovalori analitici; sorgente unica del \texttt{SparsePauliOp} riusato da VQE. & $1{\to}1,2{\to}0$\\[4pt]
\texttt{\seqsplit{vqe\_dimer.py}} & Implementazione VQE (ansatz PMA-2q$\cdot$3) per il ground state del dimero, validata contro \texttt{dimer\_exact.py}. & $1{\to}1,2{\to}0$\\[4pt]
\texttt{\seqsplit{confronto\_ansatz\_entangler.ipynb}} & Confronto sistematico di ansatz (Ry+entangler A--D, famiglie PMA) per scegliere l'ansatz canonico del dimero. & $1{\to}1,2{\to}0$\\[4pt]
\texttt{\seqsplit{vqe\_ground\_state\_test2.ipynb}} & Ground state VQE al punto di lavoro "test 2" ($b/J=0.35$, $D/J=0.80$), da riusare come stato iniziale delle correlazioni dinamiche. & $1{\to}1,2{\to}0$\\
\bottomrule
\end{longtable}

\subsection{1.2 Quantum simulation (evoluzione temporale, Suzuki--Trotter)}

\textbf{Perché Famiglia 2 — l'unica eccezione del progetto.} Il notebook
di partenza fornito con il corso
(\texttt{\seqsplit{Quantum\_simulation\_TIM\_noiseless.ipynb}}, modello di Ising
trasverso) usa già \texttt{q0,q1} per il sistema fisico e \texttt{q2} come
ancilla. Quando è stato costruito l'analogo per il dimero, la convenzione
qubit del materiale di partenza è stata mantenuta per continuità con
l'esempio fornito dal relatore — non è stata scelta ex-novo come per
\texttt{dimer\_exact.py}, è stata ereditata.

\begin{longtable}{@{}p{4.7cm}p{6.3cm}p{2.5cm}@{}}
\toprule
File & A cosa serve & Mappa\\
\midrule
\endhead
\texttt{\seqsplit{quantum\_simulation\_dimero\_teoria.md}} & Derivazione teorica completa (Pauli, fattorizzazione $H_1/H_2$, errore di Trotter $O(t^2/N)$) a monte dei due notebook. Teoria pura, nessuna convenzione qubit dichiarata. & ---\\[4pt]
\texttt{\seqsplit{quantum\_simulation\_dimero\_trotter.ipynb}} & Notebook completo: derivazione fisica del punto di lavoro R1, implementazione Trotter, verifica di convergenza. & $1{\to}0,2{\to}1$\\[4pt]
\texttt{\seqsplit{quantum\_simulation\_dimero\_trotter\_semplificato.ipynb}} & Versione ridotta (solo il richiesto dal relatore), esplorazione interattiva dei parametri con widget. & $1{\to}0,2{\to}1$\\
\bottomrule
\end{longtable}

\subsection{1.3 Correlazioni dinamiche (Hadamard test con ancilla)}

\textbf{Perché Famiglia 1, nonostante segua cronologicamente il
Trotter.} Le correlazioni dinamiche ripartono esplicitamente dalla
convenzione di \texttt{dimer\_exact.py} (stesso \texttt{SparsePauliOp},
stesso ground state VQE come stato iniziale) piuttosto che dal notebook
Trotter appena scritto: è una scelta di continuità con la parte
"statica" del progetto, non con l'ultima cosa scritta in ordine
cronologico. Per costruire il blocco $U(t)$ necessario, invece di
convertire i gate del notebook Trotter (Famiglia 2) in Famiglia 1, è
stato ricostruito da zero come esponenziale di matrice
(\texttt{scipy.linalg.expm}) — un modo per evitare il problema di
conversione piuttosto che risolverlo, dichiarato esplicitamente nel
notebook principale.

\begin{longtable}{@{}p{4.7cm}p{6.3cm}p{2.5cm}@{}}
\toprule
File & A cosa serve & Mappa\\
\midrule
\endhead
\texttt{\seqsplit{correlazioni\_dimero\_esplorazione.ipynb}} & Notebook \textbf{principale}: calcolo esplorativo di tutte le 36 combinazioni $C_{ij}^{\alpha\beta}(t)$, senza presupporre simmetrie; selezione dei correlatori per il circuito. & $1{\to}1,2{\to}0$\\[4pt]
\texttt{\seqsplit{correlazioni\_dimero\_simmetria\_U.ipynb}} & Notebook \textbf{secondario}, sincronizzato col principale: costruzione e verifica della simmetria $U=\mathrm{SWAP}\cdot(R_z(\pi)\otimes R_z(\pi))$ che dimezza le combinazioni indipendenti. & $1{\to}1,2{\to}0$\\[4pt]
\texttt{\seqsplit{verifica\_simmetrie\_correlatori.ipynb}} / \texttt{\seqsplit{\_pdf.ipynb}} & Verifica numerica indipendente (non solo simbolica) degli zeri strutturali e della simmetria $U$, a confronto con \texttt{simmetrie\_correlatori\_dimero.pdf}. & $1{\to}1,2{\to}0$\\[4pt]
\texttt{\seqsplit{circuito\_correlazioni\_tutte.ipynb}} & Notebook snello e finale: dal ground state al circuito con ancilla, misura di tutte le 36 combinazioni, visualizzazione interattiva. & $1{\to}1,2{\to}0$\\
\bottomrule
\end{longtable}

\section{Parte 2 — Trimero N=3, anello (catena chiusa)}

\textbf{Perché Famiglia 1.} Continuazione diretta della parte statica del
dimero: stessa logica di \texttt{dimer\_exact.py} applicata a $N=3$,
esplicitamente dichiarata in \texttt{trimer\_ring\_exact.py} proprio per
mantenere la corrispondenza diretta fra label Pauli e notazione fisica
$|q_1q_2q_3\rangle$ usata nella teoria (Kambe, $S_{12}^2$). Tutti e tre i
moduli VQE (\texttt{vqe\_trimer\_ring.py} e le sue due varianti)
dichiarano esplicitamente nel proprio docstring "convenzione sito
$\leftrightarrow$ qubit IDENTICA a \texttt{trimer\_ring\_exact.py}": nessuna
scelta nuova, solo eredità diretta.

\begin{longtable}{@{}p{4.7cm}p{6.3cm}p{2.5cm}@{}}
\toprule
File & A cosa serve & Mappa\\
\midrule
\endhead
\texttt{\seqsplit{teoria\_trimero\_isoscele.pdf}} & Teoria esatta (decomposizione di Kambe, spettro chiuso, campo critico, mappa dei segni). Notazione fisica, nessun codice. & $1{\to}2,2{\to}1,3{\to}0$\\[4pt]
\texttt{\seqsplit{derivazione\_stati\_kambe\_trimero.pdf}} & Derivazione esplicita degli otto autostati nella base computazionale, via operatore di abbassamento $S_-$. & $1{\to}2,2{\to}1,3{\to}0$\\[4pt]
\texttt{\seqsplit{trimer\_ring\_exact.py}} & Benchmark esatto: Hamiltoniana come \texttt{SparsePauliOp}, spettro analitico, stati di Kambe verificati contro l'operatore Qiskit, termini DM candidati. & $1{\to}2,2{\to}1,3{\to}0$\\[4pt]
\texttt{\seqsplit{vqe\_trimer\_ring.py}} & Modulo di produzione: ansatz HA e PMA con RBS e branching sullo stato iniziale. & $1{\to}2,2{\to}1,3{\to}0$\\[4pt]
\texttt{\seqsplit{vqe\_trimer\_ring\_spiegato.ipynb}} & Guida interattiva al modulo precedente. & $1{\to}2,2{\to}1,3{\to}0$\\[4pt]
\texttt{\seqsplit{vqe\_trimer\_ring\_W.py}} & Variante con il gate $W_{ij}(\theta)$ di Crippa al posto di RBS, mirror diretto di \texttt{vqe\_dimer.py}. & $1{\to}2,2{\to}1,3{\to}0$\\[4pt]
\texttt{\seqsplit{vqe\_trimer\_ring\_W\_spiegato.ipynb}} & Guida interattiva alla variante W, mirror di \texttt{vqe\_dimer\_spiegato.ipynb}. & $1{\to}2,2{\to}1,3{\to}0$\\[4pt]
\texttt{\seqsplit{vqe\_trimer\_ring\_nobranch.py}} & Ansatz PMA senza branching (stato iniziale fisso), più robusto al DM della versione a branching. & $1{\to}2,2{\to}1,3{\to}0$\\[4pt]
\texttt{\seqsplit{confronto\_ansatz\_entangler\_trimero\_unico.ipynb}} & Confronto sistematico di 10 ansatz VQE per il trimero (famiglie generiche e PMA), fidelity con e senza DM. Importa direttamente da \texttt{trimer\_ring\_exact.py}. & $1{\to}2,2{\to}1,3{\to}0$\\
\bottomrule
\end{longtable}

\textbf{Ancora da fare — quantum simulation e correlazioni dell'anello.}
A differenza del dimero, per il trimero ad anello Trotter e correlazioni
dinamiche non sono ancora stati implementati. Quando verranno svolti,
questo documento verrà aggiornato aggiungendo qui, subito sotto questa
riga, le stesse tre sottosezioni già usate per il dimero (teoria,
Trotter, correlazioni) — non in una sezione a parte in fondo al
documento, per mantenere l'ordine cronologico/tematico del lavoro.

\section{Parte 3 — Trimero N=3, catena aperta}

\textbf{Perché Famiglia 1.} Continuazione diretta dell'anello: stessa
motivazione di continuità (non reinventare la convenzione per ogni nuova
topologia) esplicitamente richiamata in
\texttt{\seqsplit{teoria\_trimero\_catena\_aperta.pdf}} e
\texttt{\seqsplit{derivazione\_stati\_kambe\_trimero\_catena.pdf}}, in modo da poter
riusare gli otto stati di Kambe già derivati senza alcuna conversione.

\begin{longtable}{@{}p{4.7cm}p{6.3cm}p{2.5cm}@{}}
\toprule
File & A cosa serve & Mappa\\
\midrule
\endhead
\texttt{\seqsplit{teoria\_trimero\_catena\_aperta.pdf}} & Teoria esatta della catena uniforme: simmetria $P_{13}$, Casimir $S_{13}^2$, spettro chiuso, campo critico $b_c=3J$, incroci veri (non anticrossing). & $1{\to}2,2{\to}1,3{\to}0$\\[4pt]
\texttt{\seqsplit{derivazione\_stati\_kambe\_trimero\_catena.pdf}} & Derivazione esplicita degli otto autostati, con gestione esplicita del riordino (coppia di classificazione $1,3$ non adiacente nel registro). & $1{\to}2,2{\to}1,3{\to}0$\\[4pt]
\texttt{\seqsplit{trimer\_chain\_exact.py}} & Benchmark esatto, analogo di \texttt{trimer\_ring\_exact.py} per la topologia a catena; 8 self-test a precisione macchina. & $1{\to}2,2{\to}1,3{\to}0$\\[4pt]
\texttt{\seqsplit{confronto\_ansatz\_entangler\_trimero\_catena.ipynb}} & Confronto sistematico di 10 ansatz VQE per la catena ($D=0$); ansatz canonico individuato (\texttt{D: Ry-RBS-Ry}). & $1{\to}2,2{\to}1,3{\to}0$\\
\bottomrule
\end{longtable}

\textbf{Ancora da fare.} VQE con termine DM (Fase 5), se richiesto dal
relatore, e — più avanti, dopo l'anello — quantum simulation e
correlazioni dinamiche sulla catena. Stessa logica di aggiornamento: le
nuove sottosezioni verranno aggiunte qui, in ordine, non in coda al
documento.

\section{Ricetta di conversione}

Per tradurre un ket o un operatore dalla Famiglia 1 alla Famiglia 2 (o
viceversa) per lo stesso numero di siti $N$: \textbf{si inverte l'ordine
dei bit/caratteri della stringa} — un'inversione globale, non uno scambio
a coppie.

\begin{itemize}
\item $N=2$: \texttt{"ab"} (Famiglia 1) $\leftrightarrow$ \texttt{"ba"}
(Famiglia 2). Es.: $|01\rangle_{F1} = |10\rangle_{F2}$ (stesso stato
fisico, stringa diversa).
\item $N=3$: \texttt{"abc"} (Famiglia 1) $\leftrightarrow$ \texttt{"cba"}
(Famiglia 2).
\end{itemize}

Per una label Pauli, stessa regola: \texttt{"XZI"} (Famiglia 1, $X$ su
sito1) diventa \texttt{"IZX"} (Famiglia 2, $X$ ancora su sito1 ma ora la
label si legge al contrario).

\textbf{Verifica pratica.} Prima di usare uno stato o un operatore
importato da un modulo con convenzione diversa: costruire
l'operatore/stato in entrambe le rappresentazioni e verificare che
l'azione su un osservabile fisico (es. $\langle S_z^{\text{totale}}\rangle$,
invariante per costruzione) coincida; se coincide solo dopo l'inversione
di stringa, la conversione è confermata corretta.

\section{Raccomandazione operativa}

Non è necessario né conveniente unificare retroattivamente le due
famiglie: sono entrambe internamente coerenti, coperte da self-test, e
riscriverle introdurrebbe rischio senza benefici a ridosso della scadenza
di settembre. La raccomandazione resta:

\begin{enumerate}
\item \textbf{Per tutto ciò che segue sulla catena aperta} (Fase 3 in
poi): restare in Famiglia 1, continuità diretta con l'anello — nessuna
conversione necessaria.
\item \textbf{Se e quando servirà collegare Trotter (Famiglia 2) con
VQE/stati di Kambe (Famiglia 1)} — rilevante per il task di quantum
simulation del dimero se lo stato iniziale verrà preparato via PMA: due
strade equivalenti già verificate praticabili, l'inversione di stringa
esplicita (sez.~5) oppure il precedente delle correlazioni dinamiche
(ricostruire $U(t)$ direttamente in Famiglia 1 come esponenziale di
matrice, senza riusare i gate del notebook Trotter così come sono).
\item \textbf{In ogni nuovo modulo}, dichiarare la convenzione scelta nel
docstring/prima cella, come già fanno tutti i file di questo documento —
è l'unica misura che ha impedito finora che questa disparità causasse un
errore reale.
\end{enumerate}

\section{Stato di verifica}

Tutti i moduli elencati sono stati verificati direttamente sul codice
sorgente o sul notebook, non per deduzione da documenti di riepilogo:
nessun punto resta aperto o "presunto". Questo documento verrà aggiornato
mano a mano che le parti mancanti (quantum simulation e correlazioni per
l'anello, VQE e seguito per la catena aperta) saranno completate,
inserendo le nuove sottosezioni nella collocazione cronologica corretta
anziché in coda.

\end{document}
